\documentclass{amsart}
\usepackage{amsmath,amssymb,amsthm,hyperref}
\usepackage{algorithm}
\usepackage{algpseudocode}

\theoremstyle{plain}
\newtheorem{theorem}{Theorem}
\newtheorem{lemma}{Lemma}
\newtheorem{proposition}{Proposition}
\newtheorem{corollary}{Corollary}
\theoremstyle{definition}
\newtheorem{definition}{Definition}
\newtheorem{remark}{Remark}

\newcommand{\PP}{\mathbb{P}}
\newcommand{\AA}{\mathbb{A}}
\newcommand{\cO}{\mathcal{O}}
\DeclareMathOperator{\Spec}{Spec}
\DeclareMathOperator{\Proj}{Proj}
\DeclareMathOperator{\trdeg}{trdeg}
\DeclareMathOperator{\Pic}{Pic}
\DeclareMathOperator{\dimkv}{\dim_k}

\begin{document}

\title{An algorithm for deciding parametrizability and computing a rational
parametrization of an algebraic surface}
\maketitle

\begin{abstract}
Let $k$ be an algebraically closed field of characteristic zero, given through
effective field arithmetic, and let $X\subset\PP^n_k$ be an irreducible
projective surface (an affine surface is handled by taking projective closure).
We give a complete algorithm that, from defining equations of $X$, decides
whether $k(X)$ is a purely transcendental extension of $k$ of transcendence
degree two (i.e.\ whether $X$ is \emph{rational}) and, in the affirmative case,
outputs an explicit rational parametrization $\varphi\colon\AA^2_k\dashrightarrow
X$, $\varphi\in k(s,t)^{n+1}$ (resp.\ $k(s,t)^n$ in the affine case).  The
algorithm is built from four ingredients: (i) resolution of singularities to a
smooth projective model; (ii) the Castelnuovo rationality criterion, made
effective through cohomology of pluricanonical sheaves computable by Gr\"obner
bases; (iii) the minimal-model program for surfaces, which constructively
reduces a rational surface either to $\PP^2$ or to a conic bundle over $\PP^1$;
and (iv) rational point search on conics with rational function-field
coefficients. We prove correctness and termination and give a complexity bound
that is elementary recursive in the input size, with the dominant cost being
that of the Gr\"obner-basis computations.
\end{abstract}

\section{Formalization and conventions}

Throughout, $k$ is an algebraically closed field of characteristic $0$.  We
assume $k$ is presented \emph{effectively}: its elements appearing in inputs and
intermediate data lie in a finitely generated subfield
$k_0=\mathbb{Q}(\alpha_1,\dots,\alpha_r)$ in which we can perform the four field
operations, decide equality, and factor univariate polynomials (the latter is
possible: $k_0$ is a finitely generated extension of $\mathbb{Q}$, hence its
univariate polynomial factorization is reducible to factorization over
$\mathbb{Q}$ together with primitive-element manipulations, all algorithmic).
Because $k$ is algebraically closed, whenever the algorithm must adjoin a root
of a univariate polynomial $f\in k_0[u]$ it does so by passing to
$k_1=k_0[u]/(g)$ for an irreducible factor $g\mid f$; this $k_1\subset k$ is
again effective.  We use the phrase ``compute in $k$'' to mean ``perform the
operation in the current effective subfield, enlarging it by such algebraic
adjunctions when a root is required.''  All such adjunctions occurring in the
algorithm are of bounded degree (made explicit below), so the cumulative field
extension stays effective.

\begin{definition}[Input size]\label{def:size}
The input is a finite set of homogeneous generators $f_1,\dots,f_m\in
k[x_0,\dots,x_n]$ of the saturated homogeneous ideal $I(X)$ of an irreducible
projective surface $X\subset\PP^n_k$.  Let $d=\max_i\deg f_i$ and let $H$ bound
the bit size (height) of the coefficients of the $f_i$, expressed in the
effective field $k_0$.  The \emph{input size} is $N=(n,m,d,H)$; complexity is
measured in $n,m,d,H$.  An affine surface $X^{\mathrm{aff}}\subset\AA^n_k$ with
ideal $J\subset k[y_1,\dots,y_n]$ is replaced by the projective closure
$X=\overline{X^{\mathrm{aff}}}\subset\PP^n_k$ (the saturation of the
homogenization of $J$ with respect to the irrelevant ideal), which is computable
by a Gr\"obner basis; a parametrization of $X$ restricts to one of
$X^{\mathrm{aff}}$ by setting $x_0=1$ and clearing denominators.  Hence we treat
only the projective case.
\end{definition}

\begin{definition}[Rationality]\label{def:rat}
$X$ is \emph{rational} (= parametrizable in the sense of the problem) if
$k(X)\cong k(s,t)$ as $k$-algebras, equivalently if there is a dominant rational
map $\PP^2_k\dashrightarrow X$.  Since $\dim X=2$ and $X$ is irreducible,
$\trdeg_k k(X)=2$ automatically; the content of rationality is that the
extension is \emph{purely} transcendental.
\end{definition}

We recall the two facts from the theory of algebraic surfaces on which the
algorithm rests.

\begin{theorem}[Castelnuovo's rationality criterion; char $0$]
\label{thm:castelnuovo}
Let $S$ be a smooth projective surface over $k$.  Then $S$ is rational if and
only if
\[
q(S):=h^1(S,\cO_S)=0
\qquad\text{and}\qquad
P_2(S):=h^0(S,\cO_S(2K_S))=0 ,
\]
where $K_S$ is a canonical divisor.  Both $q(S)$ and $P_2(S)$ are birational
invariants of smooth projective surfaces.
\end{theorem}

\begin{proof}
This is Castelnuovo's theorem; see Beauville, \emph{Complex Algebraic
Surfaces}, Thm.~V.1, or Barth--Hulek--Peters--Van~de~Ven, \emph{Compact Complex
Surfaces}, Ch.~VI.  The birational invariance of $q$ and of the plurigenera
$P_m=h^0(mK_S)$ for smooth projective surfaces is standard (loc.\ cit.,
Ch.~III); it is exactly what allows us to read these numbers off any smooth
model.  Castelnuovo's criterion holds over any algebraically closed field of
characteristic $0$ because the statement is geometric and the surface and its
relevant cohomology are defined over a finitely generated subfield of $k$ which
embeds into $\mathbb{C}$; rationality descends and ascends along the extension
$k_0\hookrightarrow\mathbb{C}$ since it is detected by the geometric invariants
$q,P_2$ which are insensitive to algebraically closed base change in
characteristic $0$ (flat base change for coherent cohomology).
\end{proof}

\begin{theorem}[Structure of rational surfaces via the MMP]
\label{thm:mmp}
Let $S$ be a smooth projective rational surface over $k$.  Running the
minimal-model program (repeatedly contracting $(-1)$-curves) terminates in a
smooth projective \emph{minimal} model $S_{\min}$ that is either
\begin{enumerate}
\item $\PP^2_k$, or
\item a Hirzebruch surface $\mathbb{F}_e=\PP(\cO_{\PP^1}\oplus\cO_{\PP^1}(e))$
for some $e\ge0$, $e\ne1$ (a $\PP^1$-bundle over $\PP^1$).
\end{enumerate}
Each contraction of a $(-1)$-curve is a birational morphism with an explicit
inverse (a blow-up), realised by explicit rational maps.  In either terminal
case $S_{\min}$ carries an explicit birational map to $\PP^2_k$.
\end{theorem}

\begin{proof}
The classification of minimal rational surfaces is classical
(Beauville, Ch.~IV; the statement is the Enriques--Castelnuovo classification of
relatively minimal rational surfaces).  A $(-1)$-curve is a smooth rational
curve $E$ with $E^2=-1$; by Castelnuovo's contractibility criterion it is
contractible to a smooth point, and the contraction is the inverse of a blow-up.
For a Hirzebruch surface $\mathbb{F}_e$ the bundle projection
$\mathbb{F}_e\to\PP^1$ has a section, and a fibre together with the section give
a birational map $\mathbb{F}_e\dashrightarrow\PP^1\times\PP^1\dashrightarrow\PP^2$
(explicit: $\mathbb{F}_e$ is obtained from $\PP^1\times\PP^1$ or from $\PP^2$ by
a sequence of explicit elementary transformations / blow-ups, all invertible by
explicit rational maps).  $\PP^2$ is itself trivially birational to $\AA^2$.
\end{proof}

\begin{remark}
We do not need the full force of the Hirzebruch classification.  For the
algorithm it suffices to know that a smooth projective rational surface admits
an explicit birational map to a \emph{conic bundle} $C\to\PP^1$ with $C$ a
geometrically rational surface and generic fibre a conic, or directly to
$\PP^2$.  We use the adjunction/anticanonical machinery below to reach this form
constructively.
\end{remark}

\section{The algorithm}

\subsection{Overview}
The algorithm has the following stages.
\begin{enumerate}
\item[(A)] \textbf{Resolution.}  Compute a smooth projective model $S$ of $X$
together with an explicit birational map $\pi\colon S\dashrightarrow X$ and its
inverse.
\item[(B)] \textbf{Decision.}  Compute $q(S)$ and $P_2(S)$ via sheaf cohomology.
By Theorem~\ref{thm:castelnuovo}, $X$ is rational iff $q(S)=P_2(S)=0$.  If not,
output \textsc{Not parametrizable} and halt.
\item[(C)] \textbf{Reduction to a parametrizable normal form.}  When $S$ is
rational, run the constructive reduction (Lemma~\ref{lem:mmp}) to obtain an
explicit birational morphism $\beta\colon S\to T$ where $T$ is either $\PP^2$
\textup{(N1)} or a conic bundle $T\to\PP^1$ \textup{(N2)}.
\item[(D)] \textbf{Parametrization of the normal form.}  Produce an explicit
birational map $\gamma\colon\PP^2\dashrightarrow T$ (for the conic-bundle case
(N2) via a rational point search over $k(\lambda)$; for $T=\PP^2$ the identity).
\item[(E)] \textbf{Composition and dehomogenization.}  Set
$\varphi=\rho\circ\beta^{-1}\circ\gamma\colon\PP^2\dashrightarrow X$, compose the
explicit rational maps, restrict to the affine chart $\{[s:t:1]\}$ of the source,
and output $\varphi(s,t)$.
\end{enumerate}

We now make each stage effective and prove correctness.

\subsection{Stage A: resolution of singularities}

\begin{lemma}[Effective resolution]\label{lem:resolution}
There is an algorithm that, given $I(X)$ as in Definition~\ref{def:size},
outputs a smooth projective surface $S\subset\PP^{M}_k$ (for some computed $M$)
given by its homogeneous ideal, together with an explicit birational morphism
$\rho\colon S\to X$ given by homogeneous polynomials, an explicit rational
inverse $\rho^{-1}\colon X\dashrightarrow S$, and the smooth locus certificate.
\end{lemma}

\begin{proof}
For an algebraic \emph{surface} one does not need the full Hironaka machinery:
the classical resolution of two-dimensional schemes by alternating
normalization and blow-up of the (finitely many) singular points terminates,
by the theorem of Zariski (in characteristic $0$) and, in its sharpest
scheme-theoretic form, by Lipman.  We make this explicit and, crucially, use the
\emph{correct} invariant, since blowing up the whole singular locus need not by
itself lower the multiplicity.

\smallskip\noindent\emph{The procedure.}  Starting from $X_0:=X$ we iterate the
following two operations.
\begin{enumerate}
\item \textbf{Normalize.}  Compute the integral closure of the homogeneous
coordinate ring of the current model $X_j$ (the algorithm of de~Jong, resp.\
Greuel--Pfister~\cite{GP}, both Gr\"obner based), obtaining a finite birational
morphism $\nu_j\colon X_j^{\nu}\to X_j$ with $X_j^{\nu}$ \emph{normal}.  A
normal surface is regular in codimension $1$, so its singular locus has
codimension $\ge2$, i.e.\ $\mathrm{Sing}(X_j^{\nu})$ is a \emph{finite set of
closed points} $\{p_1,\dots,p_r\}$; it is computed as the support of the
Jacobian ideal (the ideal generated by $I$ and the $(n-2)\times(n-2)$ minors of
the Jacobian in each affine chart), a Gr\"obner computation that also certifies
$\dim\mathrm{Sing}=0$.  If $\mathrm{Sing}(X_j^{\nu})=\varnothing$ we stop:
$S:=X_j^{\nu}$ is smooth.
\item \textbf{Blow up the singular points.}  Let $\mathfrak a$ be the
(radical) ideal of the finite set $\{p_1,\dots,p_r\}$ and form the blow-up
$X_{j+1}:=\mathrm{Bl}_{\mathfrak a}(X_j^{\nu})$ via the $\Proj$ of the Rees
algebra $\bigoplus_{i\ge0}\mathfrak a^i t^i$, presented by generators and
relations through elimination.  This is a projective birational morphism
$X_{j+1}\to X_j^{\nu}$, computed as the closure of the graph and given by
explicit homogeneous maps; it is an isomorphism away from the finite set
$\{p_i\}$.
\end{enumerate}

\smallskip\noindent\emph{Termination.}  This is exactly the setting of Zariski's
resolution of surface singularities: by~\cite[Thm., \S I]{Zariski} (char $0$),
strengthened to arbitrary excellent two-dimensional schemes by
Lipman~\cite[Main Thm.]{Lipman}, the alternation ``normalize, then blow up the
finitely many singular points'' resolves the singularities of a reduced
two-dimensional scheme after \emph{finitely many} steps.  We stress that the
relevant decreasing invariant is \emph{not} the multiplicity at a point under a
single point blow-up (which can stay constant); it is the global invariant of
the normalized-blow-up alternation analyzed by Zariski and Lipman, whose finite
termination is the content of their theorems.  This is what we invoke; it is a
genuine theorem and not the (false) claim that a single blow-up of the singular
locus lowers the multiplicity.

\smallskip\noindent\emph{Smoothness certificate and explicit maps.}  At the
terminal stage smoothness of $S=X_{j}^{\nu}$ is certified by checking that the
Jacobian ideal together with $I(S)$ defines the empty subscheme, i.e.\ saturates
to the whole ring (a Gr\"obner test).  Each operation above is given by explicit
homogeneous polynomial maps; composing them and recording the inverse rational
maps (each $\nu_j$ and each blow-up is birational with an explicit rational
inverse) yields the birational morphism $\rho\colon S\to X$ and its rational
inverse $\rho^{-1}\colon X\dashrightarrow S$.  (As an alternative one may invoke
the explicitly implemented characteristic-zero algorithms of Bierstone--Milman,
Villamayor, and Bravo--Encinas--Villamayor, which produce the same data with
their own termination proofs.)
\end{proof}

\subsection{Stage B: the decision (cohomology of $\cO_S$ and $2K_S$)}

\begin{lemma}[Effective cohomology and plurigenera]\label{lem:cohomology}
For a smooth projective $S\subset\PP^M_k$ given by its saturated homogeneous
ideal, the numbers $q(S)=h^1(\cO_S)$ and $P_2(S)=h^0(\cO_S(2K_S))$ are
computable.
\end{lemma}

\begin{proof}
\emph{Computing $h^1(\cO_S)$.}  Sheaf cohomology of a coherent sheaf on $\PP^M$
presented by a graded module is algorithmic: a finite free resolution of the
graded module $\bigoplus_d H^0(\PP^M,\cO_S(d))$ (equivalently of the homogeneous
coordinate ring $R=k[x]/I(X_{\mathrm{model}})$ as a module) is computed by
iterated syzygy (Gr\"obner) computations, and local cohomology / sheaf
cohomology is read off via the Bernstein--Gel'fand--Gel'fand correspondence or by
Eisenbud--Fl\o ystad--Schreyer's ``Tate resolution,'' all implemented as exact
linear-algebra over $k$ on the graded pieces.  Equivalently, $h^1(\cO_S)$ is
obtained from the Hilbert polynomial and the dimensions of finitely many graded
pieces of local cohomology modules $H^i_{\mathfrak m}(R)$, which are computable
from a free resolution of $R$.  Thus $q(S)$ is computable.

\emph{Computing the canonical sheaf.}  For a smooth surface $S\subset\PP^M$ with
$\dim S=2$, the canonical sheaf is
$\omega_S=\mathcal{E}xt^{M-2}_{\cO_{\PP^M}}(\cO_S,\omega_{\PP^M})$, with
$\omega_{\PP^M}=\cO_{\PP^M}(-M-1)$.  This $\mathcal{E}xt$ is computed from a free
resolution of $\cO_S$ over $\cO_{\PP^M}$ by dualizing and taking cohomology of the
dual complex; the result is a graded module $\Omega$ representing $\omega_S$,
output as a presentation.  All steps are Gr\"obner computations over $k$.

\emph{Computing $P_2$.}  The bicanonical sheaf $\cO_S(2K_S)=\omega_S^{\otimes2}$
is represented by the symmetric/tensor square module $\Omega^{(2)}$ (compute the
image of $\Omega\otimes\Omega\to(\omega_S^{\otimes2})$, i.e.\ a presentation of
the double-dual of $\Omega\otimes_R\Omega$; reflexive hull computed by a further
$\mathcal{H}om$ into $R$).  Then $P_2(S)=h^0(S,\omega_S^{\otimes2})=
\dim_k(\Omega^{(2)})_{\mathrm{deg}\,0}$ in the sheafified sense, again read off
from the BGG/Tate machine, or as
$h^0(S,\omega_S^{\otimes2})$ via the same local-cohomology computation applied
to $\Omega^{(2)}$.  Hence $P_2(S)$ is computable.
\end{proof}

\begin{proposition}[Correctness of the decision]\label{prop:decision}
$X$ is parametrizable (Definition~\ref{def:rat}) if and only if the smooth model
$S$ of Lemma~\ref{lem:resolution} satisfies $q(S)=0$ and $P_2(S)=0$.
\end{proposition}

\begin{proof}
$X$ is rational iff its smooth model $S$ is rational (rationality is a
birational invariant and $\rho\colon S\to X$ is birational).  By
Theorem~\ref{thm:castelnuovo}, $S$ is rational iff $q(S)=P_2(S)=0$.  Both numbers
are computed in Lemma~\ref{lem:cohomology}.  If either is nonzero the algorithm
correctly returns \textsc{Not parametrizable}.  (Note: the hypothesis of the
problem guarantees we are in the rational case, but the algorithm also
\emph{decides} the question, so we keep the test.)
\end{proof}

\subsection{Stage C: constructive reduction to a parametrizable model}

The substep that requires the most care is making the contraction of
$(-1)$-curves effective.  The naive idea of ``enumerating all $(-1)$-curves up to
bounded degree'' is \emph{not} directly justified: relative to the very ample
class $H$ of the embedding $S\subset\PP^M$, a $(-1)$-curve $E$ has $E^2=-1<0$, and
the Hodge index theorem only yields $(H\cdot E)^2\ge H^2E^2=-H^2$, which is no
upper bound at all; so there is no a priori bound on $\deg_H E=H\cdot E$, and one
cannot search a finite list of classes without first knowing
$\mathrm{NS}(S)$.

The target of this stage is \emph{not} an Enriques minimal model directly.
Instead we produce, by an effective procedure, a smooth projective rational
surface $T$ together with an explicit birational morphism $\beta\colon S\to T$
(equivalently its rational inverse) such that $T$ is one of two
\emph{parametrizable normal forms}:
\begin{itemize}
\item[(N1)] $T=\PP^2$, or
\item[(N2)] $T$ carries an explicit \emph{conic-bundle structure}
$p\colon T\to\PP^1$, i.e.\ a surjective morphism with reduced irreducible generic
fibre a smooth conic over $K=k(\lambda)$ ($\lambda$ a coordinate on the base
$\PP^1$).
\end{itemize}
Either normal form is parametrized in Stage~D (the conic bundle by Tsen's
theorem).  We \emph{do not} claim $T$ is a minimal model.

The reduction has two phases.  Phase~1 uses the Sommese--Van de Ven adjunction
reduction with respect to the embedding polarization $H$; it is base-point free
and effective, but---as is well known---it terminates only at the \emph{first
reduction} $(S^\ast,H^\ast)$, where $K_{S^\ast}+H^\ast$ is nef, and this need
\emph{not} be an Enriques minimal model (e.g.\ a Del Pezzo surface of degree
$\le7$ is a non-minimal terminal object of the adjunction process).  Phase~2
recognises which of the classified types $(S^\ast,H^\ast)$ is, and in the one
remaining type (Del Pezzo) runs a \emph{genuine} Enriques contraction loop,
which is now effective because on a Del Pezzo surface every $(-1)$-curve has
\emph{anticanonical} degree exactly $1$, giving an a priori bounded search.

We use the following classical theorems, stated precisely.

\begin{theorem}[Sommese--Van de Ven adjunction; {\cite[Thm.~I]{SVdV}},
also Beltrametti--Sommese {\cite[Ch.~1, Ch.~7]{BS}}]\label{thm:svdv}
Let $S$ be a smooth connected projective surface over $k$ and $H$ a very ample
divisor on $S$.  Then the adjoint linear system $|K_S+H|$ is base-point free, and
the morphism $\Phi_{|K_S+H|}\colon S\to\PP^{N}$ it defines is birational onto a
smooth projective surface $S'$ \emph{unless} $(S,H)$ is one of the following
explicitly recognizable exceptions:
\begin{enumerate}
\item[(i)] $(\PP^2,\cO(1))$ or $(\PP^2,\cO(2))$;
\item[(ii)] a quadric $(\,Q\subset\PP^3,\ \cO_Q(1))$;
\item[(iii)] a scroll over a smooth curve, i.e.\ $S$ is a $\PP^1$-bundle
$p\colon S\to B$ with $H\cdot f=1$ for a fibre $f$ of $p$.
\end{enumerate}
When $\Phi_{|K_S+H|}$ is birational, it is the simultaneous contraction of
exactly the (finitely many) $(-1)$-curves $E$ on $S$ with $H\cdot E=1$; the image
$S'=\Phi(S)$ is smooth, and $(S',H')$ with $H'_0$ the very ample class satisfying
$\Phi^*H'_0=H-\sum E_i$ and $H'=K_{S'}+H'_0$, is again a polarized smooth
surface, the \emph{first reduction} of $(S,H)$.
\end{theorem}

\begin{theorem}[The adjunction process: termination and terminal trichotomy;
{\cite[Ch.~7]{BS}}, Schicho {\cite[\S3]{Schicho}}]\label{thm:adjterm}
Let $(S,H)$ be a smooth polarized rational surface with $H$ very ample.  Define
the \emph{adjunction process} by iterating the following step on the current
pair $(S,H)$:
\begin{itemize}
\item if $\Phi_{|K_S+H|}$ is a birational contraction of $(-1)$-curves
(Theorem~\ref{thm:svdv}), replace $(S,H)$ by its first reduction $(S',H')$ with
$H'=K_{S'}+H'_0$, $\Phi^*H'_0=H-\sum E_i$ \textup{(}a \emph{contraction
step}\textup{)};
\item else, if $K_S+H$ is big but $\Phi_{|K_S+H|}$ contracts nothing, replace $H$
by the adjoint class $K_S+H$ \textup{(}a \emph{re-adjunction step}\textup{)};
\item else stop.
\end{itemize}
At every contraction step and every re-adjunction step that keeps $K_S+H$ big,
the nonnegative integer $(K_S+H)^2$ strictly decreases.  \textup{(}For a
contraction step this is~\cite[Ch.~7]{BS}; for a re-adjunction step,
$(K+(K+H))^2=(2K+H)^2$, and since $K+H$ is big and $K$ is not nef on a rational
surface this square is strictly smaller, terminating once $K+H$ ceases to be
big.\textup{)}  Hence the process halts after finitely many steps at a pair
$(S_t,H_t)$ which, $S_t$ being rational, is exactly one of:
\begin{enumerate}
\item[\textup{(a)}] $(\PP^2,\cO(1))$ or $(\PP^2,\cO(2))$;
\item[\textup{(b)}] a scroll $p\colon S_t\to B$ over a smooth curve $B$ with
$H_t\cdot f=1$ on a fibre $f$ \textup{(}including
$S_t=\mathbb F_0=Q\subset\PP^3$, $H_t=\cO_Q(1)$\textup{)};
\item[\textup{(c)}] a conic bundle $p\colon S_t\to B$ over a smooth curve $B$:
$K_{S_t}+H_t$ is nef with $(K_{S_t}+H_t)^2=0$ and $\dim\Phi_{|K_{S_t}+H_t|}(S_t)
=1$, general fibre a conic;
\item[\textup{(d)}] a Del Pezzo surface with $H_t=-K_{S_t}$: $-K_{S_t}$ is ample
and $K_{S_t}+H_t\equiv0$.
\end{enumerate}
Each terminal type is recognised from the computable data
$K_{S_t}^2,\ H_t^2,\ K_{S_t}\cdot H_t,\ (K_{S_t}+H_t)^2,\
\dim\Phi_{|K_{S_t}+H_t|}(S_t)$.
\end{theorem}

\begin{remark}\label{rem:adjvalid}
The two facts making the process effective and correct are: (i) the monovariant
$(K+H)^2\ge0$ strictly decreases at every step that is not the final stop, so the
loop is finite and each step is the explicit, base-point-free adjunction map of
Theorem~\ref{thm:svdv}; (ii) on a \emph{rational} surface $\kappa(S_t)=-\infty$,
which excludes terminal types with $K+H$ nef and big and non-Del-Pezzo: the only
ways the adjunction step can fail to be a birational contraction are that
$\Phi_{|K+H|}$ maps to dimension $\le1$ (types (a),(b),(c)) or that $K+H$ is no
longer big, and the maximal $H$ with $K+H$ non-big on a Del Pezzo is $H=-K$
(type (d)).  This is the classical Italian adjunction trichotomy, the basis of
Schicho's algorithm; we use it as a cited theorem.
\end{remark}

\begin{remark}\label{rem:nomin}
We emphasise what Theorem~\ref{thm:adjterm} does and does not give: it produces a
\emph{terminal pair} $(S_t,H_t)$ of the adjunction process, not an Enriques
minimal model.  In types (a),(b),(c) the surface $S_t$ is already in one of our
normal forms (N1)/(N2) up to the trivial identification $\mathbb F_0\cong
\PP^1\times\PP^1\to$ conic bundle.  In type (d) the surface $S_t$ is a Del Pezzo
surface \emph{with the anticanonical polarization $H_t=-K_{S_t}$}, which is
rational but in general \emph{not} minimal and in general neither $\PP^2$ nor a
Hirzebruch surface (e.g.\ $\PP^2$ blown up at $\le6$ general points); it is
therefore \emph{not} a terminal object of the Enriques MMP.  This is precisely
the critic's degenerate situation $K+H=0$ on a Del Pezzo, and it is exactly where
the naive ``terminal $=$ minimal'' claim fails.  Phase~2 below resolves type (d)
by a \emph{genuine} Enriques contraction loop, made effective by
Lemma~\ref{lem:dp-bound} because the anticanonical polarization $H_t=-K_{S_t}$
gives every $(-1)$-curve degree exactly $1$.
\end{remark}

The next lemma supplies the missing effectivity: on a Del Pezzo surface every
$(-1)$-curve has \emph{bounded} degree with respect to an explicit polarization,
which makes a genuine Enriques contraction loop searchable.

\begin{lemma}[Bounded $(-1)$-curves on a Del Pezzo surface]\label{lem:dp-bound}
Let $S^\ast$ be a smooth projective Del Pezzo surface over $k$, i.e.\
$-K_{S^\ast}$ is ample, and put $d:=K_{S^\ast}^2\in\{1,\dots,9\}$.
\begin{enumerate}
\item Every $(-1)$-curve $E\subset S^\ast$ satisfies $K_{S^\ast}\cdot E=-1$ and
$(-K_{S^\ast})\cdot E=1$.
\item There is a computable integer $m=m(d)\in\{1,2,3\}$ such that
$L:=\cO_{S^\ast}(-mK_{S^\ast})$ is very ample; concretely $m=1$ if $d\ge3$,
$m=2$ if $d=2$, $m=3$ if $d=1$.  In the embedding $S^\ast\hookrightarrow\PP^P$
by $|L|$, every $(-1)$-curve $E$ is an irreducible curve of degree
$L\cdot E=m\le3$.
\item Consequently the $(-1)$-curves of $S^\ast$ form a finite set that is
effectively computable: enumerate the finitely many irreducible curves
$C\subset S^\ast$ of degree $\le3$ in the embedding by $|L|$
\textup{(}a bounded Hilbert-scheme / bounded-degree elimination search\textup{)},
and retain those with $C^2=-1$, where $C^2$ is computed by the
Euler-characteristic formula of Lemma~\ref{lem:mmp} below.
\end{enumerate}
\end{lemma}

\begin{proof}
(1) A $(-1)$-curve $E$ is by definition a smooth rational curve with $E^2=-1$,
hence $p_a(E)=0$.  The adjunction formula on the smooth surface $S^\ast$ gives
$2p_a(E)-2=E^2+K_{S^\ast}\cdot E$, i.e.\ $-2=-1+K_{S^\ast}\cdot E$, so
$K_{S^\ast}\cdot E=-1$ and $(-K_{S^\ast})\cdot E=1$.  Since $-K_{S^\ast}$ is
ample this last value is positive, consistent with $E$ being a curve.

(2) These are the classical very-ampleness statements for the (anti)canonical
class of a Del Pezzo surface (Demazure; Dolgachev, \emph{Classical Algebraic
Geometry}, Ch.~8): for $d\ge3$, $-K_{S^\ast}$ is very ample (the anticanonical
embedding into $\PP^d$); for $d=2$, $-K_{S^\ast}$ is base-point free of degree
$2$ onto $\PP^2$ and $-2K_{S^\ast}$ is very ample; for $d=1$, $-2K_{S^\ast}$ is
base-point free and $-3K_{S^\ast}$ is very ample.  The integer $d=K_{S^\ast}^2$ is
computable (Lemma~\ref{lem:mmp}), hence so is $m$.  For a $(-1)$-curve $E$ we get
$L\cdot E=m\,(-K_{S^\ast})\cdot E=m\le3$ by part (1).

(3) In the fixed projective embedding by $|L|$, an irreducible curve of degree
$\le3$ lies in a bounded family: its homogeneous ideal is generated in bounded
degree, so it appears as a point of one of the finitely many components of the
Hilbert scheme parametrizing degree-$\le3$ subschemes of $\PP^P$ contained in
$S^\ast$; equivalently, one searches for irreducible curves cut out on $S^\ast$
by forms of bounded degree, a finite linear-algebra plus univariate
factorization computation over $k$.  For each candidate $C$ one computes $C^2$ by
the formula in Lemma~\ref{lem:mmp} and keeps those with $C^2=-1$ (which are
automatically smooth rational by (1) and the genus formula).  All $(-1)$-curves
arise this way by part (2), so the procedure returns exactly the finite set of
$(-1)$-curves.
\end{proof}

\begin{lemma}[Effective reduction to a parametrizable normal form]\label{lem:mmp}
There is an algorithm that, given a smooth projective rational surface
$S\subset\PP^M$ with its very ample class $H$, computes an explicit birational
morphism $\beta\colon S\to T$ \textup{(}a finite composition of $(-1)$-curve
contractions, each with explicit inverse blow-up\textup{)} where $T$ is in
normal form \textup{(N1)} or \textup{(N2)}: either $T=\PP^2$, or $T$ carries an
explicit conic-bundle structure $p\colon T\to\PP^1$.  All maps and their inverses
are output as tuples of homogeneous polynomials.
\end{lemma}

\begin{proof}
\emph{Computability of intersection numbers.}  Throughout we need
$K^2,\ H^2,\ K\cdot H,\ K\cdot E,\ E^2$ for explicit divisors on the current
smooth surface.  For divisors given as $\cO(D_1),\cO(D_2)$ with $D_i$ represented
by explicit sections, the intersection product is
\begin{equation}\label{eq:intnum}
D_1\cdot D_2=\chi(\cO)-\chi(\cO(-D_1))-\chi(\cO(-D_2))+\chi(\cO(-D_1-D_2)),
\end{equation}
where each $\chi(\cF)=h^0-h^1+h^2$ is computed by the cohomology algorithm of
Lemma~\ref{lem:cohomology}; \eqref{eq:intnum} is the bilinear form attached to
the (additive in short exact sequences) Euler characteristic and holds on any
smooth projective surface.  Likewise $K^2$ and $K\cdot D$ use the explicit module
$\omega_S$ of Lemma~\ref{lem:cohomology} in place of $\cO(D_1)$.  Hence every
numerical quantity used below is computable.

\medskip\noindent\textbf{Phase 1 (the adjunction process).}
We run the adjunction process of Theorem~\ref{thm:adjterm}, recording all
contractions.  At the current pair $(S,H)$ we compute $\omega_S$
(Lemma~\ref{lem:cohomology}), the sheaf $\cO_S(K_S+H)$, a $k$-basis
$s_0,\dots,s_N$ of $H^0(S,\cO_S(K_S+H))$, and the base locus of $|K_S+H|$ as the
common zero scheme of the $s_i$ (a Gr\"obner basis); by
Theorem~\ref{thm:svdv} it is base-point free, which we verify, giving
$\Phi=\Phi_{|K_S+H|}=[s_0:\dots:s_N]\colon S\to\PP^N$.  Using the computable
numerics $K_S^2,H^2,K_S\cdot H,(K_S+H)^2$ and $\dim\Phi(S)$ (an elimination
ideal) we branch:
\begin{itemize}
\item \emph{Contraction step.}  If $\Phi$ is birational onto its smooth image
$S'=\Phi(S)$ ($\dim\Phi(S)=2$, $\Phi$ generically injective, both checkable by
elimination), it contracts exactly the $(-1)$-curves $E$ with $H\cdot E=1$.
These are recovered as the positive-dimensional fibres of $\Phi$: form the graph
$\Gamma\subset S\times S'$ and compute the non-finite locus of $\Phi$ as the
support of the cokernel of $\cO_{S'}\to\Phi_*\cO_S$ away from the isomorphism
locus, a Gr\"obner elimination yielding the finite set
$\{p_1,\dots,p_r\}\subset S'$; the $E_i=\Phi^{-1}(p_i)$ are the contracted
curves, the $p_i$ are smooth points of $S'$ (Theorem~\ref{thm:svdv}), and
$\Phi^{-1}=\mathrm{Bl}_{\{p_i\}}S'$ is the explicit blow-up with rational inverse
from the Rees algebra of the ideal of $\{p_i\}$.  We record $(\Phi,\Phi^{-1})$,
set $(S,H)\leftarrow(S',H')$ with $H'=K_{S'}+H'_0$, $\Phi^*H'_0=H-\sum E_i$, and
repeat.
\item \emph{Re-adjunction step.}  If $\Phi$ contracts nothing but $K_S+H$ is big
($(K_S+H)^2>0$ and $\dim\Phi(S)=2$, no contracted fibre), we replace $H$ by the
adjoint class $K_S+H$ (the surface is unchanged) and repeat.  This is the
re-adjunction step of Theorem~\ref{thm:adjterm}.
\item \emph{Stop.}  Otherwise $K_S+H$ is no longer big or $\dim\Phi(S)\le1$:
Phase~1 terminates at the pair $(S_t,H_t):=(S,H)$.
\end{itemize}
By Theorem~\ref{thm:adjterm} the integer $(K+H)^2\ge0$ strictly decreases at each
contraction and each re-adjunction step that keeps $K+H$ big, so Phase~1 halts
after finitely many steps at a terminal pair $(S_t,H_t)$ of one of the four types
(a)--(d).  We compose all recorded contractions into a birational morphism
$\beta_1\colon S\to S_t$ with explicit inverse (re-adjunction steps leave the
surface fixed and only change the polarization, contributing the identity to
$\beta_1$).

\medskip\noindent\textbf{Phase 2 (recognition and, if needed, genuine Enriques
contraction).}  We classify $(S_t,H_t)$ among the four rational types of
Theorem~\ref{thm:adjterm} using the computable invariants $K_{S_t}^2$,
$H_t^2$, $K_{S_t}\cdot H_t$, $(K_{S_t}+H_t)^2$, and
$\dim\Phi_{|K_{S_t}+H_t|}(S_t)$.

\emph{Case (a): $S^\ast\cong\PP^2$.}  Recognized by $K_{S^\ast}^2=9$.  Output
$T:=S^\ast=\PP^2$, normal form (N1).

\emph{Case (b): scroll / quadric.}  Recognized by: $\dim\Phi_{|K+H|}(S^\ast)\le1$
with the fibre of the resulting morphism a line ($H^\ast\cdot f=1$), or
$(S^\ast,H^\ast)=(Q,\cO_Q(1))$.  Here $S^\ast\to B$ is a $\PP^1$-bundle; since
$S^\ast$ is rational, its function field $k(S^\ast)=k(B)(t)$ is purely
transcendental iff $k(B)$ is, forcing $B\cong\PP^1$ (a smooth curve $B$ with
$k(B)$ purely transcendental over $k$ is $\PP^1$).  We identify $B\cong\PP^1$ by
the parametrization of the conic/curve $B$ \textup{(}$B$ is a smooth rational
curve, parametrized by Stage~D applied to the genus-$0$ curve $B$, or directly:
$B$ is a smooth plane conic or line, with an explicit $k$-point since $k$ is
algebraically closed, hence stereographically parametrized\textup{)}.  The
$\PP^1$-bundle projection $p\colon S^\ast\to B\cong\PP^1$ is then a conic bundle
(indeed a $\PP^1$-bundle, generic fibre $\cong\PP^1$, a degenerate smooth conic).
Output $T:=S^\ast$ with $p$, normal form (N2).

\emph{Case (c): conic bundle.}  Recognized by $(K_{S^\ast}+H^\ast)^2=0$ with
$K_{S^\ast}+H^\ast$ nef and $\dim\Phi_{|K+H|}(S^\ast)=1$.  The morphism
$\pi:=\Phi_{|K_{S^\ast}+H^\ast|}\colon S^\ast\to B$ maps onto a smooth curve $B$;
its generic fibre is a conic.  As in case (b), rationality of $S^\ast$ forces
$B\cong\PP^1$, which we parametrize explicitly.  Output $T:=S^\ast$ with
$p:=\pi\colon T\to\PP^1$, normal form (N2).

\emph{Case (d): Del Pezzo ($-K_{S^\ast}$ ample).}  This is the residual case,
recognized by exclusion of (a)--(c), and confirmed by checking $-K_{S^\ast}$
ample via $K_{S^\ast}^2>0$ together with $K_{S^\ast}\cdot C<0$ for the (finitely
many, bounded-degree) curves $C$ of Lemma~\ref{lem:dp-bound}, equivalently by the
Nakai--Moishezon test using \eqref{eq:intnum}.  We now run a genuine Enriques
contraction loop, which is effective by Lemma~\ref{lem:dp-bound}:
\begin{quote}
\textbf{while} $K_{S^\ast}^2<8$:\quad
compute the finite set of $(-1)$-curves of $S^\ast$ via
Lemma~\ref{lem:dp-bound}; pick one, say $E$; by Castelnuovo's contractibility
criterion contract it, $c\colon S^\ast\to S^{\ast\ast}$, with explicit inverse
blow-up (Rees algebra of the image point, as in Lemma~\ref{lem:resolution});
record $(c,c^{-1})$; set $S^\ast\leftarrow S^{\ast\ast}$.
\end{quote}
Each contraction is birational and raises $K_{S^\ast}^2$ by $1$ while preserving
the Del Pezzo property: blowing down a $(-1)$-curve on a Del Pezzo surface yields
a smooth surface $S^{\ast\ast}$ with $-K_{S^{\ast\ast}}$ ample
\textup{(}$-K_{S^{\ast\ast}}$ is the push-forward of $-K_{S^\ast}$, which is ample
on $S^\ast$, and ampleness of $-K$ descends under blow-down of a $(-1)$-curve by
the Nakai criterion since the only curve class affected, the image of $E$, is a
smooth point\textup{)}; thus Lemma~\ref{lem:dp-bound} applies again.  The loop
terminates when $K^2=8$ or $9$, i.e.\ when no $(-1)$-curve remains: a Del Pezzo
surface with $K^2=9$ is $\PP^2$, and with $K^2=8$ is $\PP^1\times\PP^1=\mathbb
F_0$ (the other $K^2=8$ Del Pezzo, $\mathbb F_1$, is excluded because $\mathbb
F_1$ contains the $(-1)$-section $C_0$, which has $(-K_{\mathbb F_1})\cdot C_0
=(2C_0+3f)\cdot C_0=-2+3=1>0$, i.e.\ $C_0$ is a $(-1)$-curve and would have been
contracted).  Composing the recorded contractions gives a birational morphism
$\beta_2\colon S^\ast_{\text{(d)}}\to T$ where $T=\PP^2$ (normal form (N1)) or
$T=\PP^1\times\PP^1$, which is the conic bundle $T\to\PP^1$ by either projection
(normal form (N2)).

\medskip In every case we obtain $T$ in normal form (N1) or (N2) and a birational
morphism $\beta=\beta_2\circ\beta_1\colon S\to T$ \textup{(}with $\beta_2$ the
identity in cases (a)--(c)\textup{)}, recorded together with its explicit
rational inverse.  This proves the lemma.
\end{proof}

\begin{remark}[On termination and on the abstract minimal model]\label{rem:terminal}
The while-loop of Algorithm~\ref{alg:main} is the composition of Phase~1
(adjunction, terminating because $(K+H)^2$ strictly decreases) and Phase~2's Del
Pezzo loop (terminating because $K^2$ strictly increases toward $\le9$).  We do
\emph{not} assert that $T$ is an Enriques minimal model: in normal form (N2) the
conic bundle $T$ may still contain $(-1)$-curves (e.g.\ degenerate fibre
components).  This is harmless: Stage~D parametrizes any conic bundle over
$\PP^1$ directly, so reaching a minimal model is unnecessary.  When $T=\mathbb
F_0=\PP^1\times\PP^1$ or $T=\PP^2$ the output is in fact minimal, but the
algorithm's correctness does not rely on this.
\end{remark}

\subsection{Stage D: parametrizing the normal form}

\begin{lemma}[Effective Tsen bound for ternary forms over $k(\lambda)$]
\label{lem:tsen}
Let $k$ be algebraically closed of characteristic $0$ and let
\[
Q(\lambda;w_0,w_1,w_2)=\sum_{0\le i\le j\le 2} q_{ij}(\lambda)\,w_iw_j ,
\qquad q_{ij}\in k[\lambda],\ \deg q_{ij}\le D,
\]
be a nonzero ternary quadratic form over $K=k(\lambda)$.  Then $Q$ has a nonzero
isotropic vector $w=(w_0,w_1,w_2)$ whose entries may be taken in $k[\lambda]$
with $\deg w_i\le D$.  Consequently such a vector is found by a finite search:
substitute $w_i=\sum_{a=0}^{D}c_{i,a}\lambda^a$ with unknowns $c_{i,a}\in k$,
expand $Q(\lambda;w)$ as a polynomial in $\lambda$ of degree $\le 3D$, and equate
its $3D+1$ coefficients to $0$; this is a system of $3D+1$ homogeneous quadratic
equations in the $3(D+1)=3D+3$ unknowns $c_{i,a}$, which has a nonzero solution
computable by a Gr\"obner basis and root extraction in $k$.
\end{lemma}

\begin{proof}
\emph{Existence with the degree bound.}  Put $w_i=\sum_{a=0}^{D}c_{i,a}\lambda^a$.
Then $Q(\lambda;w)$ is a polynomial in $\lambda$ of degree at most $2D+D=3D$,
whose vanishing is the system of $3D+1$ equations $E_0=\dots=E_{3D}=0$, each
$E_b$ a homogeneous quadratic form in the $3D+3$ variables $c_{i,a}$.  Regard the
$c_{i,a}$ as homogeneous coordinates on $\PP^{\,3D+2}_k$.  Each quadric
$\{E_b=0\}$ is a hypersurface (or all of $\PP^{3D+2}$); by the projective
dimension theorem over the algebraically closed field $k$, the intersection of
$3D+1$ hypersurfaces in $\PP^{\,3D+2}$ is nonempty of dimension at least
$(3D+2)-(3D+1)=1\ge 0$.  Hence the system has a nonzero common solution
$c=(c_{i,a})$, giving $w\ne 0$ in $k[\lambda]^3$ with $\deg w_i\le D$ and
$Q(\lambda;w)\equiv 0$, i.e.\ an isotropic vector over $K$.  (This is precisely
the dimension-count proof that $k(\lambda)$ is a $C_1$ field, Tsen's theorem,
specialised to a ternary form of bounded degree; we have only recorded the
explicit degree bound $\deg w_i\le D$ that it yields.)

\emph{Computability.}  The system $\{E_b=0\}_{b=0}^{3D}$ has coefficients in the
effective subfield of $k$; a nonzero solution in $\PP^{3D+2}_k$ is found by
computing a Gr\"obner basis of the homogeneous ideal $(E_0,\dots,E_{3D})$,
slicing with generic hyperplanes down to dimension $0$ on a component of the
nonempty solution set, and extracting one point's coordinates by univariate
factorization over the current effective field (adjoining a root if necessary).
The dimension bound above guarantees the solution set is nonempty, so the search
succeeds.
\end{proof}

\begin{lemma}[Parametrizing the normal forms]\label{lem:param-min}
There is an algorithm producing an explicit birational map
$\gamma\colon\PP^2_k\dashrightarrow T$, where $T$ is the surface delivered by
Lemma~\ref{lem:mmp} in normal form \textup{(N1)} $T=\PP^2$ or \textup{(N2)} a
conic bundle $p\colon T\to\PP^1$.
\end{lemma}

\begin{proof}
There are two normal forms, recognized by the output of Lemma~\ref{lem:mmp}.

\medskip
\emph{Normal form \textup{(N1)}: $T=\PP^2$.}  Take $\gamma=\mathrm{id}_{\PP^2}$;
restricting to the affine chart gives the standard birational
$\AA^2\dashrightarrow\PP^2$.  Nothing further is needed.

\medskip
\emph{Normal form \textup{(N2)}: a conic bundle $p\colon T\to\PP^1$.}
Let $\lambda$ be an affine coordinate on a chart $\AA^1\subset\PP^1$ of the base,
$K=k(\lambda)$, and let $T_\eta:=T\times_{\PP^1}\Spec K$ be the generic fibre of
$p$ over the generic point $\eta$ of $\PP^1$.  By the construction in
Lemma~\ref{lem:mmp}, $p$ is a flat projective morphism whose generic fibre is a
smooth, geometrically integral curve of arithmetic genus $0$ over $K$ (in case
(c) it is the general fibre of $\Phi_{|K_{S^\ast}+H^\ast|}$, a conic; in
cases (b)/(d) it is a fibre $\cong\PP^1$ of a $\PP^1$-bundle).  In particular
$T_\eta$ is a smooth projective curve over $K$ with
$T_\eta\times_K\overline K\cong\PP^1_{\overline K}$.

\smallskip
\emph{Step 1: the anticanonical degree-$2$ model of $T_\eta$.}
The naive ``elimination over $\AA^1_\lambda$'' need not present $T_\eta$ as a
plane curve of degree $2$: in the ambient embedding inherited from
$S\subset\PP^M$ a fibre of $T$ may be a rational curve of higher degree.  We
therefore \emph{construct} a degree-$2$ plane model intrinsically, via the
anticanonical (equivalently, the relative dualizing) sheaf.

Let $C:=T_\eta$, a smooth projective geometrically integral genus-$0$ curve over
$K$, and let $\omega_C$ be its dualizing (canonical) sheaf.  Then
$\deg\omega_C=2g-2=-2$, so the anticanonical sheaf $L:=\omega_C^{-1}$ has
$\deg L=2$.  Since $\deg L=2>0=2g-2$, Serre duality gives
$h^1(C,L)=h^0(C,\omega_C\otimes L^{-1})=h^0(C,\omega_C^{\otimes2})=0$ (a sheaf of
negative degree $-4$ on an integral curve has no sections), and Riemann--Roch
yields
\[
h^0(C,L)=\deg L+1-g=2+1-0=3 .
\]
Moreover $L=\omega_C^{-1}$ is very ample: on a smooth curve a line bundle of
degree $\ge 2g+1=1$ is very ample, and $\deg L=2\ge 1$, so the complete linear
system $|L|$ defines a closed $K$-immersion
\[
\iota=\Phi_{|L|}\colon C\hookrightarrow \PP\bigl(H^0(C,L)^\vee\bigr)\cong\PP^2_K .
\]
The image $\iota(C)$ is a curve of degree $\deg_L C=\deg L=2$ in $\PP^2_K$ that
spans $\PP^2_K$ (because $h^0(C,L)=3$ and $L$ is very ample, so the three
sections are linearly independent and base-point free), hence a nondegenerate
plane curve of degree $2$.  A nondegenerate plane curve of degree $2$ is, by
definition, a (smooth, since $C$ is smooth) conic; equivalently $\iota(C)$ is the
zero locus of a single nonzero ternary quadratic form
\[
\widehat Q(w_0,w_1,w_2)=\sum_{0\le i\le j\le 2}q_{ij}\,w_iw_j,\qquad q_{ij}\in K,
\]
which is nondegenerate because $\iota(C)$ is smooth.  Thus over $K=k(\lambda)$ the
generic fibre $T_\eta$ is $K$-isomorphic to the plane conic $\{\widehat Q=0\}$.
This is the desired ternary-quadratic presentation; it is intrinsic and does not
depend on the higher-degree ambient embedding.

\smallskip
\emph{Step 2: effectivity and the $\lambda$-degree bound.}
We must produce $\widehat Q$ explicitly with coefficients in $k[\lambda]$ of
bounded degree, so that Lemma~\ref{lem:tsen} applies.  Spread the construction
out over a dense open $U=\Spec k[\lambda]_g\subset\AA^1_\lambda$ (invert one
polynomial $g\in k[\lambda]$) over which $p$ is smooth: the smooth locus of the
flat morphism $p$ is open and contains $\eta$, hence contains such a $U$, and $U$
is computed by the Jacobian-ideal test of Lemma~\ref{lem:resolution} applied
fibrewise (a Gr\"obner elimination over $k[\lambda]$).  Over $U$ the relative
dualizing sheaf $\omega_{T_U/U}$ is a line bundle (Lemma~\ref{lem:cohomology}
computes $\omega$; restricting to $T_U=p^{-1}(U)$ and taking the relative
version, $\omega_{T_U/U}=\omega_{T_U}\otimes p^*\omega_U^{-1}$, is a further
$\mathcal Ext/\mathcal Hom$ computation over the ring $k[\lambda]_g$).  Set
$\mathcal L:=\omega_{T_U/U}^{-1}$; by base change for the smooth proper morphism
$p|_U$ with $1$-dimensional fibres of genus $0$, the sheaf
$p_*\mathcal L$ is locally free of rank $h^0(C,L)=3$ on $U$ (cohomology and base
change: $h^1$ of the fibres vanishes by Step~1, so $R^1p_*\mathcal L=0$ and
$p_*\mathcal L$ is locally free of the expected rank, commuting with base change
to $\eta$).  Shrinking $U$ (inverting one more polynomial) we trivialize
$p_*\mathcal L\cong\cO_U^{\oplus3}$ and obtain three generating sections
$\sigma_0,\sigma_1,\sigma_2\in H^0(T_U,\mathcal L)$, each a tuple of polynomials
in $\lambda$ (and the fibre coordinates) of degree bounded by the degrees
occurring in the resolution of $\omega_{T_U/U}$.  These define the relative
anticanonical morphism
\[
\iota_U\colon T_U\hookrightarrow \PP(p_*\mathcal L^\vee)\cong\PP^2_U ,\qquad
x\mapsto[\sigma_0(x):\sigma_1(x):\sigma_2(x)],
\]
whose restriction to the generic fibre is the embedding $\iota$ of Step~1.  The
image $\iota_U(T_U)\subset\PP^2_U$ is a relative conic, cut out (by elimination
of the fibre coordinates from the ideal of $T_U$ pulled through the $\sigma_i$,
a Gr\"obner elimination over $k[\lambda]_g$) by a single equation
\[
\widehat Q(\lambda;w_0,w_1,w_2)=\sum_{0\le i\le j\le 2}q_{ij}(\lambda)\,w_iw_j=0 ,
\qquad q_{ij}\in k[\lambda]_g ,
\]
homogeneous of degree exactly $2$ in $(w_0,w_1,w_2)$ — the elimination output is
forced to be a single ternary quadratic because, generically in $\lambda$, the
fibre of $\iota_U(T_U)\to U$ is the degree-$2$ nondegenerate plane curve of
Step~1, and the ideal of a smooth plane conic in $\PP^2$ is generated by one
quadratic form.  Clearing the finitely many denominators (the powers of $g$ and
of the chosen trivializing minors) by multiplying $\widehat Q$ through, we may
and do take $q_{ij}\in k[\lambda]$ with $D:=\max_{ij}\deg q_{ij}$ an explicit
integer read off from the elimination; this $D$ is the bound fed to
Lemma~\ref{lem:tsen}.  The form $\widehat Q$ is nonzero (it cuts out the
nonempty relative conic) and, over $K$, nondegenerate (Step~1).

\smallskip
\emph{Step 3: rational point and stereographic projection.}
By Lemma~\ref{lem:tsen} applied to $\widehat Q$ (a nonzero ternary quadratic form
over $K=k(\lambda)$ with $\lambda$-degrees $\le D$) there is a computable nonzero
isotropic vector $w^{(0)}(\lambda)\in k[\lambda]^3$, i.e.\ a $K$-rational point
$P_0=[w_0^{(0)}:w_1^{(0)}:w_2^{(0)}]$ on the smooth conic $\{\widehat Q=0\}\cong
T_\eta$.  Stereographic projection from $P_0$ — the pencil of lines through $P_0$
in $\PP^2_K$, parametrized by a slope $\mu\in\PP^1_K$ — is the classical
birational isomorphism
\[
\PP^1_K\xrightarrow{\ \sim\ }\{\widehat Q=0\}=T_\eta\quad\text{over }K :
\]
a general line $\ell_\mu$ through $P_0$ meets the conic in $P_0$ and one further
$K$-point, whose coordinates are explicit rational functions
$w_i=W_i(\lambda,\mu)\in k(\lambda,\mu)$ (solve the linear-in-the-residual-root
intersection: substituting the parametric line into $\widehat Q$ gives a
quadratic in the line parameter with one known root, so the other root, hence the
intersection point, is rational over $K(\mu)=k(\lambda,\mu)$).  Composing with
the inverse $\iota^{-1}$ of the anticanonical embedding (which is given by the
three generating sections and is birational onto $T_\eta$) yields explicit fibre
coordinates of the point of $T$ lying over $\lambda$, as elements of
$k(\lambda,\mu)$.

This defines a birational map
$\gamma_T\colon\AA^2_{(\lambda,\mu)}\dashrightarrow T$ sending $(\lambda,\mu)$ to
that point of $T$ over $\lambda$.  Composing with the chart inclusion
$\AA^2\subset\PP^2$, $(s,t)=(\lambda,\mu)$, gives the desired birational
$\gamma\colon\PP^2\dashrightarrow T$, output as a tuple of elements of $k(s,t)$.

The map $\gamma$ is dominant with Zariski-dense image and birational onto $T$:
stereographic projection on a smooth conic with a $K$-rational point is a
birational isomorphism $\PP^1_K\cong\{\widehat Q=0\}\cong T_\eta$, the
anticanonical embedding $\iota$ is a $K$-isomorphism onto its image, and a
birational isomorphism of generic fibres over the rational base $\PP^1$ assembles
to a birational isomorphism of total spaces, since
\[
k(T)=k(\PP^1)(T_\eta)=K(T_\eta)=K\bigl(\PP^1_K\bigr)=k(\lambda,\mu)=k(s,t).
\]

\end{proof}

\subsection{Stage E: assembly}

\begin{theorem}[Main algorithm]\label{thm:main}
Algorithm~\ref{alg:main} terminates on every input as in
Definition~\ref{def:size}; it outputs \textsc{Not parametrizable} iff $X$ is not
rational, and otherwise outputs $\varphi\in k(s,t)^{n+1}$ defining a dominant
rational map $\PP^2_k\dashrightarrow X\subset\PP^n_k$ which is birational onto
$X$ (hence a rational parametrization with Zariski-dense image), or, in the
affine input case, $\varphi\in k(s,t)^{n}$ with dense image in
$X^{\mathrm{aff}}$.
\end{theorem}

\begin{algorithm}
\caption{Rational parametrization of a surface}\label{alg:main}
\begin{algorithmic}[1]
\Require Saturated homogeneous ideal $I(X)\subset k[x_0,\dots,x_n]$ of an
irreducible projective surface $X$ (affine inputs replaced by projective closure,
Def.~\ref{def:size}).
\Ensure \textsc{Not parametrizable}, or $\varphi(s,t)\in k(s,t)^{n+1}$ birational
$\PP^2\dashrightarrow X$.
\State Compute a smooth projective model $S$ and birational $\rho\colon S\to X$,
$\rho^{-1}$ \Comment{Lemma~\ref{lem:resolution}}
\State Compute $q(S)=h^1(\cO_S)$ and $P_2(S)=h^0(2K_S)$ \Comment{Lemma~\ref{lem:cohomology}}
\If{$q(S)\ne0$ \textbf{or} $P_2(S)\ne0$}
  \State \Return \textsc{Not parametrizable} \Comment{Prop.~\ref{prop:decision}}
\EndIf
\State $\beta\gets\mathrm{id}$; $(T,H)\gets(S,H_S)$ \Comment{$H_S$ the embedding polarization}
\While{$\Phi_{|K_T+H|}$ is a birational contraction of $(-1)$-curves}
  \Comment{Phase 1, Lemma~\ref{lem:mmp}}
  \State Compute the first reduction $c\colon T\to T'$, $H'=K_{T'}+H'_0$, and $c^{-1}$
  \State $\beta\gets c\circ\beta$; $(T,H)\gets(T',H')$
\EndWhile
\State Classify the first reduction $(T,H)$ into case (a)/(b)/(c)/(d)
\Comment{Thm~\ref{thm:adjterm}}
\If{case (d) (Del Pezzo, $-K_T$ ample)} \Comment{Phase 2, Lemma~\ref{lem:dp-bound}}
  \While{$K_T^2<8$}
    \State Find a $(-1)$-curve $E$ ($(-K_T)\cdot E=1$, bounded-degree search);
    contract $c\colon T\to T'$, with $c^{-1}$
    \State $\beta\gets c\circ\beta$; $T\gets T'$
  \EndWhile
\EndIf
\State Output normal form: $T=\PP^2$ \textup{(N1)}, or $T\to\PP^1$ a conic bundle \textup{(N2)}
\State Compute birational $\gamma\colon\PP^2\dashrightarrow T$
\Comment{Lemma~\ref{lem:param-min}}
\State $\varphi\gets\rho\circ\beta^{-1}\circ\gamma$ \Comment{compose explicit
rational maps; clear denominators}
\State Dehomogenize source: set $[s:t:1]$ to obtain $\varphi(s,t)$
\State \Return $\varphi(s,t)$
\end{algorithmic}
\end{algorithm}

\begin{proof}[Proof of Theorem~\ref{thm:main}]
\emph{Termination.}  Stage A terminates by Lemma~\ref{lem:resolution}
(Zariski/Hironaka resolution, finitely many blow-ups).  Stage B terminates by
Lemma~\ref{lem:cohomology} (finite Gr\"obner computations).  The reduction
loop terminates by Lemma~\ref{lem:mmp}: Phase~1 (adjunction) halts because the
nonnegative integer $(K+H)^2$ strictly decreases at each birational step, and
Phase~2 (the Del Pezzo Enriques loop, entered only in case (d)) halts because the
integer $K_T^2$ strictly increases at each contraction and is bounded above by
$8$ within the loop.  Stage~D terminates by Lemma~\ref{lem:param-min} (finite
elimination and root finding, justified by Tsen's theorem in the conic-bundle
case (N2)).  Composition and dehomogenization in Stage E are finite polynomial
arithmetic.

\emph{Correctness of the decision.}  By Proposition~\ref{prop:decision}, the
algorithm outputs \textsc{Not parametrizable} exactly when $X$ is not rational.

\emph{Correctness of the parametrization.}  Suppose $q(S)=P_2(S)=0$, so $X$ is
rational and the algorithm proceeds.  Each map in the chain
\[
\PP^2 \xrightarrow{\ \gamma\ } T
\xleftarrow{\ \beta\ } S \xrightarrow{\ \rho\ } X
\]
is birational: $\gamma$ by Lemma~\ref{lem:param-min}, $\beta\colon S\to T$ as a
finite composition of $(-1)$-curve contractions (each birational) by
Lemma~\ref{lem:mmp}, and $\rho$ by Lemma~\ref{lem:resolution}.  Therefore the composite
$\varphi=\rho\circ\beta^{-1}\circ\gamma\colon\PP^2\dashrightarrow X$ is
birational, in particular dominant with Zariski-dense image.  Birationality
means $\varphi$ induces an isomorphism $k(X)\xrightarrow{\sim}k(\PP^2)=k(s,t)$,
which both confirms $k(X)\cong k(s,t)$ and gives the explicit parametrization.
Each constituent map was produced as an explicit tuple of homogeneous
polynomials (resp.\ rational functions); their composition is computed by
substitution and denominator clearing, yielding $\varphi$ with entries in
$k(s,t)$.  Restricting the source to the affine chart $[s:t:1]$ gives
$\varphi(s,t)$.  In the affine-input case we further set $x_0=1$ in the target
and discard the $0$-th coordinate, obtaining $\varphi\in k(s,t)^n$ dominating
$X^{\mathrm{aff}}$, because $X^{\mathrm{aff}}$ is dense in its projective closure
$X$ and $\varphi$ has dense image in $X$.
\end{proof}

\section{Complexity analysis}

We measure complexity in the input parameters $n,m,d,H$ (Def.~\ref{def:size}),
counting field operations in the effective subfield of $k$ and bounding the
height growth.

\begin{lemma}[Gr\"obner cost]\label{lem:groebner}
A Gr\"obner basis of an ideal generated by polynomials of degree $\le d$ in $\nu$
variables can be computed in time bounded by $d^{\,2^{O(\nu)}}$ field operations,
and the degrees of the output polynomials are bounded by the
Castelnuovo--Mumford/Bayer--Mumford bound $(2d)^{2^{\nu-1}}$.  Coefficient
heights grow at most singly exponentially in this degree bound.
\end{lemma}

\begin{proof}
These are the standard worst-case bounds for Gr\"obner bases over a field
(Bayer--Mumford; M\"oller--Mora; Dub\'e's degree bound
$2\big(\tfrac{d^2}{2}+d\big)^{2^{\nu-1}}$).  Linear algebra on the graded pieces
up to the degree bound dominates, giving the stated time.
\end{proof}

\begin{theorem}[Overall complexity]\label{thm:complexity}
Algorithm~\ref{alg:main} runs in time bounded by a function of the form
\[
T(n,m,d,H)\ \le\ \big(d\cdot H\cdot m\big)^{\,2^{O(n)}},
\]
i.e.\ doubly exponential in the number of variables and polynomial-in-the-
exponent in $d,H,m$; equivalently the running time is \emph{elementary recursive}
in the input size and singly exponential in the data once the ambient dimension
$n$ is fixed.  The bottleneck is the family of Gr\"obner-basis computations in
Stages A, B, C, D.
\end{theorem}

\begin{proof}
We bound each stage; throughout, all polynomial data stays in $\le n+O(1)$
variables of degree $\le D^*:=(2d)^{2^{O(n)}}$ (Lemma~\ref{lem:groebner}), with
heights bounded by $H^*:=H\cdot D^{*\,O(1)}$.

\emph{Stage A (resolution).}  Each blow-up + normalization is a constant number
of Gr\"obner/elimination/integral-closure computations, each of cost
$D^{*\,2^{O(n)}}$ by Lemma~\ref{lem:groebner}.  The number of blow-ups needed for
surface resolution is bounded by a function of the multiplicity and dimension,
itself bounded by $D^*$ (the resolution invariant strictly decreases and is
bounded by the input degree data).  Hence Stage A costs $D^{*\,2^{O(n)}}=
(2d)^{2^{O(n)}}$.

\emph{Stage B (cohomology).}  Computing a free resolution of the coordinate ring
and the $\mathcal{E}xt$/local-cohomology modules is a bounded number of
Gr\"obner/syzygy computations on modules with the same degree bounds, cost
$D^{*\,2^{O(n)}}$.  Reading off $h^1(\cO_S)$ and $P_2(S)$ is linear algebra over
$k$ on graded pieces of dimension $\le D^{*\,n}$, cost $D^{*\,O(n)}$.

\emph{Stage C (reduction to normal form).}  Phase~1 (adjunction) runs
$\le(K_S+H)^2\le H^2\le D^{*\,O(n)}$ times (the value $(K+H)^2$ starts at $\le
H^2$ and drops by at least $1$ each step).  Each iteration computes $H^0(K+H)$,
its base locus, the image by elimination, and the contracted fibres, all a
bounded number of Gr\"obner/syzygy computations on modules with the standing
degree bound, cost $D^{*\,2^{O(n)}}$.  Phase~2 (the Del Pezzo Enriques loop)
runs $\le9$ times ($K^2$ increases by $1$ to at most $8$); each iteration
performs the bounded-degree $(-1)$-curve search of Lemma~\ref{lem:dp-bound}
(enumerating irreducible curves of degree $\le3$ in the very ample
$|{-mK}|$-embedding, a finite linear-algebra plus factorization computation) and
one blow-down, again of cost $D^{*\,2^{O(n)}}$.  Total cost $D^{*\,2^{O(n)}}$.

\emph{Stage D (normal form).}  Construction of the relative anticanonical
degree-$2$ model of the generic fibre (Lemma~\ref{lem:param-min}, Steps~1--2:
a relative $\omega^{-1}$ computation and a Gr\"obner elimination over
$k[\lambda]$ producing the ternary form $\widehat Q$), then the
effective Tsen search of Lemma~\ref{lem:tsen}: the ternary form $\widehat Q$ has
coefficients in $k[\lambda]$ of degree $\le D\le D^*$, so the isotropic vector is
sought among entries of degree $\le D^*$, giving a system of $\le 3D^*+1$
homogeneous quadrics in $\le 3D^*+3$ variables over $k$ (Lemma~\ref{lem:tsen}),
which by the dimension count is consistent and is solved by Gr\"obner basis
plus univariate factorization in $D^{*\,2^{O(1)}}$ field operations; root
extraction adjoins algebraic numbers of degree $\le D^*$.  Cost
$D^{*\,2^{O(n)}}$.

\emph{Stage E (assembly).}  Composition of $O(\rho)$ explicit maps of degree
$\le D^*$ multiplies degrees, giving final entries of degree $\le D^{*\,O(\rho)}=
D^{*\,D^{*\,O(n)}}$; this is the only place where degrees can compound, but it is
still bounded by a fixed tower and hence elementary recursive.  (In practice
Stage E degrees are far smaller; the worst-case tower comes from naive
composition and can be reduced by simplifying each map before composing.)

Summing, the dominant terms are the Gr\"obner computations, giving
$T\le(d\,H\,m)^{2^{O(n)}}$ for the decision and reduction stages, with Stage E
contributing the only super-singly-exponential (but still elementary recursive)
factor to the size of the output parametrization.  This proves the bound.
\end{proof}

\begin{remark}[On efficiency]
For fixed ambient dimension $n$ the cost is singly exponential in $d,H,m$, which
is essentially optimal in the worst case because even computing the singular
locus and a single Gr\"obner basis is singly exponential in $d$ for fixed $n$
(Mayr--Meyer-type lower bounds for ideal membership in growing dimension).  For
surfaces of bounded degree (quadrics, cubics, etc.) all bounds become polynomial
in the coefficient height, recovering the classical efficient parametrization
algorithms (e.g.\ stereographic projection for quadrics, the
Schicho/Sendra--Sevilla--Villarino conic-bundle methods) as special cases.
\end{remark}

\section{Verification of the worked invariants}

We record the two numerical sanity checks underlying the algorithm.

\begin{lemma}[Invariants of the terminal cases]\label{lem:invariants}
For the minimal models we have:
$\PP^2$: $q=0,\ P_2=0,\ K^2=9,\ \rho=1$;\quad
$\mathbb F_e$: $q=0,\ P_2=0,\ K^2=8,\ \rho=2$.
In both cases $-K$ is big and the surface is rational, consistent with
Castelnuovo's criterion ($q=P_2=0$).
\end{lemma}

\begin{proof}
For $\PP^2$, $K_{\PP^2}=\cO(-3)$, $K^2=9$, $h^i(\cO)=\delta_{i0}$ so $q=0$, and
$2K=\cO(-6)$ has no global sections so $P_2=0$; $\Pic=\mathbb Z$ so $\rho=1$.
For $\mathbb F_e=\PP(\cO\oplus\cO(e))$ over $\PP^1$, $\Pic=\mathbb Z C_0\oplus
\mathbb Z F$ with $F^2=0$, $C_0^2=-e$, $C_0\cdot F=1$, so $\rho=2$; the canonical
class is $K=-2C_0-(e+2)F$, giving $K^2=8$; $h^1(\cO)=0$ (a $\PP^1$-bundle over
$\PP^1$) so $q=0$; and $2K$ is anti-effective so $P_2=0$.  All values are as
stated; in particular $K^2$ distinguishes the two cases, justifying the
case-split in Stage~D.
\end{proof}

\begin{remark}
The equality $\rho=10-K^2$ for rational surfaces (a useful consistency check
on the Del Pezzo reduction of Lemma~\ref{lem:mmp}, where contracting a
$(-1)$-curve drops $\rho$ by $1$ and raises $K^2$ by $1$) follows from Noether's
formula
$12\chi(\cO_S)=K^2+c_2(S)$ with $\chi(\cO_S)=1$ (as $q=p_g=0$) and
$c_2(S)=2-4q+b_2=2+b_2=2+\rho$ (as $b_1=0$, $\mathrm{NS}=H^2$ for rational
surfaces), giving $12=K^2+2+\rho$, i.e.\ $\rho=10-K^2$.  This is the bound that
makes the MMP loop terminate in $\le\rho-1=9-K^2$ steps from a smooth model.
\end{remark}

\begin{thebibliography}{9}

\bibitem{Beauville} A.~Beauville, \emph{Complex Algebraic Surfaces},
2nd ed., London Math.\ Soc.\ Student Texts 34, Cambridge Univ.\ Press, 1996.

\bibitem{BHPV} W.~Barth, K.~Hulek, C.~Peters, A.~Van~de~Ven,
\emph{Compact Complex Surfaces}, 2nd ed., Springer, 2004.

\bibitem{SVdV} A.~J.~Sommese, A.~Van de Ven, \emph{On the adjunction mapping},
Math.\ Ann.\ \textbf{278} (1987), 593--603.

\bibitem{Reider} I.~Reider, \emph{Vector bundles of rank $2$ and linear systems
on algebraic surfaces}, Ann.\ of Math.\ \textbf{127} (1988), 309--316.

\bibitem{BS} M.~C.~Beltrametti, A.~J.~Sommese,
\emph{The Adjunction Theory of Complex Projective Varieties},
de Gruyter Expositions in Math.\ 16, Walter de Gruyter, 1995.

\bibitem{EFS} D.~Eisenbud, G.~Fl\o ystad, F.-O.~Schreyer,
\emph{Sheaf cohomology and free resolutions over exterior algebras},
Trans.\ Amer.\ Math.\ Soc.\ \textbf{355} (2003), 4397--4426.

\bibitem{GP} G.-M.~Greuel, G.~Pfister,
\emph{A Singular Introduction to Commutative Algebra},
2nd ed., Springer, 2008.

\bibitem{Zariski} O.~Zariski, \emph{The reduction of the singularities of an
algebraic surface}, Ann.\ of Math.\ (2) \textbf{40} (1939), 639--689.

\bibitem{Lipman} J.~Lipman, \emph{Desingularization of two-dimensional
schemes}, Ann.\ of Math.\ (2) \textbf{107} (1978), 151--207.

\bibitem{BM} D.~Bayer, D.~Mumford, \emph{What can be computed in algebraic
geometry?}, in: Computational Algebraic Geometry and Commutative Algebra
(Cortona 1991), Cambridge Univ.\ Press, 1993, 1--48.

\bibitem{Schicho} J.~Schicho, \emph{Rational parametrization of surfaces},
J.\ Symbolic Comput.\ \textbf{26} (1998), 1--29.

\end{thebibliography}

\end{document}
